% No preamble. Prefix fqi-. Cross-references only to this fragment and fcc-.
\section{Exact finite local quantum induction}
\label{fqi-section}

This section combines the actual common-domain finite macro integral,
finite-total quotient descent, the fixed-outer-radius force of
Notebook30, and Theorem~\ref{fcc-theorem}. The result concerns the
selected homogeneous \emph{local} macrostep and its specified
holomorphic continuation. It does not identify the full Wilson
pushforward, other field sectors, or arbitrary mixed histories.

\subsection{Finite geometry, totals and the actual map}
\label{fqi-setup-section}

Fix the primitive ratio $\ell$, an elementary composite ratio
$R=\ell^s$, auxiliary grouping $J$, and exponent $a_0=1/16$
satisfying the strict common-domain budgets. At stage $j$ let
\begin{equation}
 M_j=M_0R^j,\qquad Q_j=Q_0/R^j.
 \label{fqi-scales}
\end{equation}
The preceding homogeneous classical depth is $M_j$; the cubic
current-unit torus period is $Q_j$, after resetting spacing to one.
All periods are compatible integers and $J\mid Q_j$. Constants
may depend on the fixed geometry and color conventions, but are
uniform in depth and in the permitted periods above a lower
threshold.

On the period-$Q$ torus use the finite activity space and quotient
of \eqref{fcc-norm-total}. For clarity, the raw and normalized totals
are
\begin{equation}
 E^{\rm raw}(A)=\sum_HE_H(A|_H),\qquad
 \Phi_E(A)=E^{\rm raw}(A)-E^{\rm raw}(0),\qquad
 \|E\|_{Q;a_0,b}=\sup_i\sum_{H\ni i}e^{a_0|H|}\|E_H\|_b.
 \label{fqi-totals}
\end{equation}
There are finitely many coordinates and connected $J$-cell carriers
at each regulator. Equality of $\Phi_E$ on the full product ball
is a closed linear kernel, so the quotient is Banach. We impose
local gauge, translation, hypercubic and global color identities
on totals, not individual coefficients. Use a simple Lie algebra
for the scalar Maxwell functional; a finite componentwise version
handles a semisimple product.

Let $S_{M,Q}$ be the actual finite homogeneous ordinary minimum
action, with $S_{M,Q}(0)=0$. It is not replaced by a canonical
periodization of an independently selected infinite reference.
The actual common-domain construction supplies
\begin{align}
 \mathcal F_j(E)&=\mathcal T_jE+\mathcal G_j(E),\notag\\
 \mathcal G_j(E)&=\Delta_j+\mathcal Q_{{\rm det},j}
                                      +\mathcal L_j(Q_E).
 \label{fqi-map}
\end{align}
Here $\mathcal T_j$ is pre-gauge substitution, $Q_E$ is the exact
inherited fluctuation difference, and $\mathcal L_j$ is the
strict Gaussian logarithm with the fixed actual base exponent,
precision, retained window and chart. The coupling $g_j$ and
threshold $T_{{\rm cut},j}$ are held fixed when comparing inputs
to this one map. Field-independent factors are tracked separately.

On the original real common overlap, this is the exact chosen
local macro integral. On the common radius-$b$ ball it is its
proved holomorphic continuation, evaluated through pre-gauge
fields inside the old activity domain. Identification throughout
that ball with a globally defined gauge-orbit density requires
the extra orbit-representation hypothesis of the common-domain
theorem. It is not inferred from a quotient identity.

\subsection{Exact finite quotient descent}
\label{fqi-descent-section}

\begin{lemma}[Finite descent of the actual quantum update]
\label{fqi-descent-lemma}
The actual finite maps $\mathcal F_j$, $\mathcal T_j$ and
$\mathcal G_j$ descend to the invariant finite quotient on
their admissible ball. No infinite periodization, kernel limit
or translation averaging is required.
\end{lemma}
\begin{proof}
If $[E]=[E']$, then on the entire input product ball
\begin{equation}
 E^{\rm raw}(A)-E'^{\rm raw}(A)=c,
 \qquad c=E^{\rm raw}(0)-E'^{\rm raw}(0).
 \label{fqi-constant-difference}
\end{equation}
The actual pre-gauge field $a_p(w)$ and all fluctuation fields
$f_\lambda(w,\xi)$ used in the inherited construction lie in
that ball on the fixed retained window and the established
interpolation domains. Consequently
\begin{equation}
 E^{\rm raw}(f_1)-E^{\rm raw}(f_0)
       =E'^{\rm raw}(f_1)-E'^{\rm raw}(f_0).
 \label{fqi-same-insertion}
\end{equation}
The scalar inherited Gaussian insertion is identical, while
the exterior background terms differ by $c$. The covariance,
base action, chart, sharp window and measure factors are unchanged.
The exact finite integrals therefore differ only by $e^c$.
Their normalized holomorphic logarithms agree on the connected
domain, using the common nonvanishing branch. A raw logarithm
constant belongs to the vacuum normalization and disappears
from the quotient. The localized arrays need not agree.

Classical substitution descends immediately from
\eqref{fqi-constant-difference} evaluated at $a_p(w)$;
subtracting it gives descent of $\mathcal G_j$.
The exact real-sector symmetry and holomorphic Ward continuation
preserve the stipulated total identities. No carrierwise gauge
identity is assumed.
\end{proof}

The constant in \eqref{fqi-constant-difference} must be compensated
in a raw density, as specified below. For a real class, a chosen
representative can be made real without increasing its norm by
averaging $E_H(w)$ with $\overline{E_H(\overline w)}$.

\subsection{Uniform finite difference and forcing estimates}
\label{fqi-estimates-section}

Let $\bar\epsilon$ be the family ceiling of the strict two-input
comparison. For quotient inputs $x,y$ of norm at most
$\bar\epsilon/4$, choose near-minimal representatives $Y$ of $y$
and $D$ of $x-y$, and put $X=Y+D$. For arbitrarily small
$\delta>0$,
\begin{align}
 \|Y\|&\le\bar\epsilon/4+\delta,\notag\\
 \|D\|&\le\|x-y\|+\delta\le\bar\epsilon/2+\delta,\notag\\
 \|X\|&\le3\bar\epsilon/4+2\delta<\bar\epsilon.
 \label{fqi-compatible-representatives}
\end{align}
Both complete inputs are admissible and have invariant totals.
Covariance of their arrays is unnecessary at finite volume.

The actual Gaussian interpolation retains the base exponent and
precision. In particular impose
\begin{equation}
 v\le\frac{\eta_*}{64e},\qquad
 a_Q\bar\epsilon\le\frac{\eta_*}{128e},\qquad
 a_Q=16C_TC_R\frac{g_jT_{{\rm cut},j}}b.
 \label{fqi-comparison-budgets}
\end{equation}
The inherited difference is linear in $D$. The finite
Banach-valued logarithm Cauchy estimate gives
\begin{equation}
 \|\mathcal G_j(x)-\mathcal G_j(y)\|_{Q_{j+1}}
 \le L_{\mathrm{quant},j}\|x-y\|_{Q_j},\qquad
 L_{\mathrm{quant},j}
   =C_{\mathrm{quant}}\frac{g_jT_{{\rm cut},j}}b,\quad
 C_{\mathrm{quant}}=1024eS_1C_TC_R.
 \label{fqi-quantum-Lipschitz}
\end{equation}
Apply the family bound to the compatible $X,Y$, use
Lemma~\ref{fqi-descent-lemma}, and let $\delta$ decrease to zero.
No Fatou or local-limit argument is needed at a fixed regulator.
The external determinants cancel in this difference, and no
equal-low-jet hypothesis is used.

For the zero input, choose one outer field radius $\rho$ after
$R,J$ with all strict classical and Gaussian margins. Fix the
same actual threshold, charts, covariance and branch before
restricting to $b$. Notebook30's finite outer-domain force has
a uniform anchored bound $H_*$ on $\rho$. Its total linear
jet vanishes by semisimple adjoint invariance. Removing
coefficientwise degrees zero and one and using homogeneous
phase projection gives, for $b\le\rho/2$,
\begin{equation}
 \|\mathcal G_j(0)\|_{Q_{j+1}}\le C_Fb^2,
 \qquad C_F=2H_*/\rho^2.
 \label{fqi-force}
\end{equation}
This is an exactly equivalent normalized representative, not
a low-jet approximation. Its vacuum constant is kept separately.
The full force need not be $O(g^2)$: the pivot and Gaussian
determinant terms can persist at zero coupling.

\subsection{Finite references and the classical block}
\label{fqi-classical-section}

At stage $j$, use the finite Maxwell ratio of
\eqref{fcc-projection}, normalized against the actual reference
$S_j=[S_{M_j,Q_j}]$. For periods above its fixed threshold,
\begin{equation}
 \beta_j(S_j)=1,\qquad
 \|\beta_j\|\le C_\beta b^{-2},\qquad
 P_j=I-S_j\beta_j,\qquad
 \|P_j\|\le p:=1+C_\beta C_S,\qquad
 \|S_j\|\le C_Sb^2.
 \label{fqi-projections}
\end{equation}
The $b^{-2}$ factor is retained and cancels only in the
combined projection estimate. All constants are uniform in
the permitted periods and preceding depth.

The actual finite two-minimum comparison and ordinary
composition give
\begin{equation}
 \mathcal T_jS_j=S_{j+1},\qquad
 P_{j+1}\mathcal T_jP_j=P_{j+1}\mathcal T_j.
 \label{fqi-reference-telescope}
\end{equation}
In dimension four the fine action has the identical
$\eta^{d-4}=1$ normalization in both unit conventions.
The identity concerns the full actual action, not merely
its quadratic symbol. Put $A_j=P_{j+1}\mathcal T_j$.
The finite classical semigroup for invariant totals yields
\begin{equation}
 A_{j+q-1}\cdots A_jP_j
     =P_{j+q}\mathcal T_{M_j,R^q,Q_{j+q}}P_j.
 \label{fqi-block-telescope}
\end{equation}
The direct map on the right is classical substitution.
It does not replace the intervening quantum integrals or
their sharp windows by a one-shot integral.

Write $L=R^q$. Theorem~\ref{fcc-theorem} gives the block bound
\begin{equation}
 \theta(q,N):=p^2\left[
 C_{\mathrm{low}}(J)L^{-\alpha}+\frac{6C_{\mathrm{high}}}L
 +C_{\mathrm{err}}(J)(L^2+L+1)e^{-c(J)LN}\right],
 \qquad N=Q_{j+q}.
 \label{fqi-classical-block-bound}
\end{equation}
It retains the graded finite Ward defects and wrapping tails
at the same $a_0,b$. Fix $J$, choose a finite $q$, and then
a compatible $N_0$ above all reference, collar and period
thresholds so that
\begin{equation}
                 \theta(q,N)\le1/4\qquad(N\ge N_0).
 \label{fqi-period-threshold}
\end{equation}
Only afterward choose $b$ within the direct $L$ requirements
of \eqref{fcc-radius} and all common domains. One-step
transport has uniform bounds
\begin{equation}
 \|\mathcal T_j\|\le B_T,\qquad
 \|A_j\|\le B,\qquad B\ge\max(1,pB_T).
 \label{fqi-one-step-bound}
\end{equation}
The finite projection telescope remains exact although the
period and functional $\beta_j$ change at each stage.
It requires no identity $\beta_{j+1}\mathcal T_j=\beta_j$.

\subsection{A uniform finite block envelope}
\label{fqi-envelope-section}

Define the extracted maps
\begin{equation}
 \Psi_j(x)=P_{j+1}\mathcal F_j(x)=A_jx+H_j(x),\qquad
 H_j=P_{j+1}\mathcal G_j.
 \label{fqi-extracted-map}
\end{equation}
Let $\nu$ bound $pL_{\mathrm{quant},j}$ over the step data.
Equations \eqref{fqi-quantum-Lipschitz}--\eqref{fqi-force} give
\begin{equation}
 \operatorname{Lip}H_j\le\nu,\qquad
 \|H_j(0)\|\le C_0b^2,\qquad C_0=pC_F.
 \label{fqi-perturbation-bounds}
\end{equation}
Set
\begin{equation}
 K=B+1,\qquad L_q=\sum_{l=0}^{q-1}K^l,\qquad
 K_{\mathrm{amp}}=\max\{1,2C_0L_q\},\qquad
 C_{\mathrm{int}}=K^qK_{\mathrm{amp}}+C_0L_q.
 \label{fqi-envelope-constants}
\end{equation}
Choose the common $b$ so that $C_{\mathrm{int}}b^2<\bar\epsilon/4$,
and finally restrict the weak-coupling range so every strict
budget holds and
\begin{equation}
 \nu\le\min\{1,(4qK^{q-1})^{-1}\}.
 \label{fqi-nu-choice}
\end{equation}
For example, a sufficiently large fixed multiple of
$\sqrt{\log(1/g)}$ with a fixed lower threshold has
$gT(g)/b\to0$, $gT(g)^3\to0$ and $\tau(T(g))\to0$.
The outer force remains uniform on this smaller coupling range.

\begin{theorem}[Finite local induction with exact small representatives]
\label{fqi-theorem}
Use the actual finite maps, references and strict domains above.
Choose $q,N_0,b$ and the weak-coupling range in the stated order.
For every compatible finite history of $n$ steps with
\begin{equation}
 Q_0=R^nQ_n,\qquad Q_n\ge N_0,
 \label{fqi-history-periods}
\end{equation}
and initial extracted class satisfying
\begin{equation}
 \beta_0(x_0)=0,\qquad
 \|x_0\|\le r_b:=K_{\mathrm{amp}}b^2,
 \label{fqi-initial-ball}
\end{equation}
the following conclusions hold for couplings in the permitted
weak range.

Every full $q$-step block preserves the radius-$r_b$ extracted
ball and has Lipschitz constant at most $1/2$ when compared
with the same prescribed coupling data. Every intermediate
class has norm at most $C_{\mathrm{int}}b^2<\bar\epsilon/4$.
At every stage one can choose an exactly equivalent connected
finite representative on the same $a_0,b$ with norm at most
$2C_{\mathrm{int}}b^2<\bar\epsilon/2$. Thus every actual next
input and Gaussian budget remains valid. The changes of raw
constant are compensated by the exact vacuum factors below.

The constants are independent of the permitted number of
blocks and of preceding depth. The result describes the
selected homogeneous local germ and its continuation; its
global measure interpretation retains the stated extra
orbit-domain requirements.
\end{theorem}

\begin{proof}
For a block starting at stage $j$ with an input in
\eqref{fqi-initial-ball}, the one-step bound gives, inductively
for $0\le k\le q$,
\begin{equation}
 \|x_{j+k}\|
 \le K^kK_{\mathrm{amp}}b^2+C_0b^2\sum_{l<k}K^l
 \le C_{\mathrm{int}}b^2<\bar\epsilon/4.
 \label{fqi-intermediate-bound}
\end{equation}
This verifies the next map's domain before that map is invoked.
It also applies to the zero-started comparison trajectory.

For two trajectories using the same couplings, variation of
constants bounds their difference by
\begin{align}
 \|\delta_q\|
 &\le\left[\theta(q,Q_{j+q})
       +\nu\sum_{k=0}^{q-1}B^{q-1-k}(B+\nu)^k\right]
                                             \|\delta_0\|\notag\\
 &=\left[\theta(q,Q_{j+q})+(B+\nu)^q-B^q\right]
                                             \|\delta_0\|\notag\\
 &\le\left[\frac14+q\nu K^{q-1}\right]\|\delta_0\|
 \le\frac12\|\delta_0\|.
 \label{fqi-Duhamel}
\end{align}
Only the initial full classical product uses
\eqref{fqi-classical-block-bound}; every shorter product
uses powers of $B$. The block output at zero has norm at
most $C_0L_qb^2\le r_b/2$. Comparing with that output
therefore proves invariance of the radius-$r_b$ ball.
The output is extracted at every step by definition.

Repeat by full blocks; any final partial block satisfies
\eqref{fqi-intermediate-bound}. Condition
\eqref{fqi-history-periods} ensures every full block output
has the lower period required by
\eqref{fqi-period-threshold}. On each fixed starting torus
there are finitely many permitted steps. Uniformity as that
torus grows does not assert an infinite sequence on one
fixed finite regulator.

For representative selection, choose at stage $j$ a finite
family $E_j$ within $C_{\mathrm{int}}b^2$ of the quotient infimum.
Then
\begin{equation}
 \|E_j\|_{Q_j;a_0,b}
        \le2C_{\mathrm{int}}b^2<\bar\epsilon/2.
 \label{fqi-chosen-family}
\end{equation}
The infimum need not be attained. This family is holomorphic
on the full same $b$ ball, has connected carriers at the
same exponent, and has exactly the required normalized total.
Its total symmetries are unchanged. A real class permits the
norm-nonincreasing real averaging described earlier.

Use $E_j$ as the next input. Its norm is strictly inside the
actual family ceiling, so the inherited support estimates
and Gaussian theorem apply. Lemma~\ref{fqi-descent-lemma}
shows that this choice changes no next quotient output.
After the next step and actual finite reference subtraction,
choose another such family. Every generated correction remains
in the total. Since equality is exact, the positive
near-minimal tolerance creates no accumulating approximation
error and need not form a summable sequence. The raw constant
is handled by \eqref{fqi-vacuum-update} below.
\end{proof}

This selection occurs at each finite regulator; it is not
an inference from a small infinite class. Every winding
contribution is part of the finite total estimated by
Theorem~\ref{fcc-theorem} and the finite quantum bounds.
The assertion does not say that an untouched canonical
cluster array remains small. Selection is existential in
a finite Banach quotient, not a proposed algorithm for
finding a minimizing family. Exact equality and the uniform
norm are the properties needed for this induction.

The envelope also applies along an adaptively updated physical
coupling sequence for as long as it stays weak. For the bound
on a single such history one may freeze its realized coupling
values and use the auxiliary zero-started comparison with those
same values. This does not give a difference contraction
between histories that have different updated couplings.

\subsection{Exact vacuum compensation and coupling bookkeeping}
\label{fqi-vacuum-section}

Write the finite input density on the common chart as
\begin{equation}
 Z_j\exp[-\gamma_jS_j+E_j^{\rm raw}],\qquad
 \gamma_j=g_j^{-2}.
 \label{fqi-input-density}
\end{equation}
The finite exact chosen local step, with its known scalar
normalization $N_j$, produces
\begin{equation}
 Z_jN_j\exp[-\gamma_jS_{j+1}+F_j^{\rm raw}].
 \label{fqi-output-density}
\end{equation}
Let
\[
 \beta_{\rm out}=\beta_{j+1}([F_j^{\rm raw}]),\qquad
 x_{j+1}=[F_j^{\rm raw}]-\beta_{\rm out}[S_{j+1}],
\]
and choose the small family $E_{j+1}$ representing this class.
Since $S_{j+1}(0)=0$, equality of normalized totals gives
\begin{align}
 F_j^{\rm raw}(w)
   &=\beta_{\rm out}S_{j+1}(w)
                       +E_{j+1}^{\rm raw}(w)+c_j,\notag\\
 c_j&=F_j^{\rm raw}(0)-E_{j+1}^{\rm raw}(0).
 \label{fqi-constant-compensation}
\end{align}
The equality is on the entire common product domain. Hence
the exact updates preserving \eqref{fqi-output-density} are
\begin{equation}
 \gamma_{j+1}=\gamma_j-\beta_{\rm out},\qquad
 Z_{j+1}=Z_jN_j e^{c_j}.
 \label{fqi-vacuum-update}
\end{equation}
If an input representative is replaced by
$E'^{\rm raw}=E^{\rm raw}+c$ before integration, replace
$Z_j$ by $Z_je^{-c}$. Both signs are fixed by the raw
exponential, not by a convention about the quotient.
On real inputs use the compatible real logarithm branches;
constant $2\pi i$ changes in another compatible convention
do not change the exponentials. No convergence of an extensive
vacuum free energy is asserted.

If the couplings are freely prescribed without satisfying
\eqref{fqi-vacuum-update}, the theorem still bounds the
corresponding projected maps. It does not then assert
repeated factorization of one unchanged density. Along
the actual update, the amplitude bounds hold only while
the couplings stay in their proved weak range.

For the unprojected value along the envelope,
\begin{equation}
 \|\mathcal F_j(x_j)\|
 \le[B_TC_{\mathrm{int}}+C_F+L_{\mathrm{quant}}C_{\mathrm{int}}]b^2,
 \qquad L_{\mathrm{quant}}=\sup_j L_{\mathrm{quant},j}.
 \label{fqi-unprojected-bound}
\end{equation}
Separating $\mathcal T_jx_j$, $\mathcal G_j(0)$ and
$\mathcal G_j(x_j)-\mathcal G_j(0)$ proves this bound.
The finite Maxwell functional consequently gives
\begin{equation}
 |\gamma_{j+1}-\gamma_j|
 \le C_\beta[B_TC_{\mathrm{int}}+C_F
                              +L_{\mathrm{quant}}C_{\mathrm{int}}].
 \label{fqi-increment-bound}
\end{equation}
This is uniform in the finite regulator but has no sign.
For a specified finite number $n$ of steps, an initial
inverse coupling exceeding the weak threshold by $n$ times
this bound suffices for weak containment on that horizon.
It does not select a desired final coupling, establish an
averaged positive drift or control an infinite physical
trajectory. Comparing different running-coupling histories
would require an additional coupling-parameter comparison.

\subsection{Precise closure and remaining obligations}
\label{fqi-scope-section}

The theorem closes the finite \emph{representative-budget}
interface for this homogeneous local scheme. At every allowed
stage the chosen connected family has norm $O(b^2)$ at the
same carrier exponent and field radius. Every actual next
input and Gaussian budget holds; finite normalized-total
identities are exact; vacuum constants are compensated by
\eqref{fqi-constant-compensation}--\eqref{fqi-vacuum-update}.
The q-block amplitude contraction is at most one half for
equal prescribed coupling data.

This does not prove smallness of an untouched carrier array,
identify a one-shot composite sector with sequential primitive
sectors, supply large-field or mixed-history terms, or show
that the local germs represent every global Wilson history.
Those interpretations require their own orbit-domain and
measure identities. Every retained window remains precisely
the one in the common-domain theorem.

A physical density trajectory must obey
\eqref{fqi-vacuum-update} and remain weak. Its ultraviolet
tuning and beta drift, other field sectors, vacuum limits,
observables, nontriviality and a continuum measure remain
separate. None follows solely from the finite representative
induction proved here.
